% Part: incompleteness
% Chapter: theories-computability
% Section: consis-with-q

\documentclass[../../../include/open-logic-section]{subfiles}

\begin{document}

\olfileid{inc}{tcp}{con}

\olsection{$\Th{Q}$ نال کٹھے بے تضاد نظریے ناقابلِ فیصلہ نیں}

اگلے قضیے مطابق $\Th{Q}$ ناقابلِ فیصلہ اے، تے ہر اوہ نظریہ
وی ناقابلِ فیصلہ اے جیہڑا $\Th{Q}$ نال کٹھا بے تضاد ہووے۔

\begin{thm}
فرض کرو $\Th{T}$ حساب دی زبان دا کوئی نظریہ اے جیہڑا
$\Th{Q}$ نال کٹھا بے تضاد اے (یعنی $\Th{T} \cup \Th{Q}$
بے تضاد اے)۔ تاں $\Th{T}$ ناقابلِ فیصلہ اے۔
\end{thm}

\begin{proof}
یاد رکھو کہ $\Th{Q}$ دے مسلّمیاں دا سیٹ محدود اے، یعنی
$!Q_1$، \dots، $!Q_8$۔ انہاں دی جگہ اسیں اکّو مسلّمہ وی
رکھ سکدے آں، $!E = !Q_1 \land
\dots \land !Q_8$۔

فرض کرو $\Th{T}$ اک قابلِ فیصلہ نظریہ اے جیہڑا~$\Th{Q}$ نال
کٹھا بے تضاد اے۔ منّ لئیے کہ
\[
C = \Setabs{!A}{\Th{T} \Proves !E \lif !A}.
\]
اسیں وکھاواں گے کہ $C$، $\Th{Q}$ تے $\Th{\bar Q}$ نوں
وکھ کرن والا قابلِ حساب سیٹ اے۔ ایس توں تضاد نکلے گا۔
دتے ہوئے جملے توں مقرر شرطیہ بنانا مؤثر کم اے۔ فرض کیتے
فیصلہ کرن والے طریقے نال پتا لگ سکدا اے کہ ایہ شرطیہ ثابت
ہو سکدا اے یا نہیں۔ ایس لئی سیٹ قابلِ حساب اے۔ پہلاں جے
$!A$، $\Th{Q}$ وچ اے، تاں $!A$، $\Th{Q}$ دے مسلّمیاں توں
ثابت ہو سکدا اے۔ اخذ دے قضیے مطابق، پہلے درجے دی منطق وچ
$!E \lif !A$ دا !!a{derivation} اے۔ ایس لئی $!A$، $C$ وچ اے۔

دوجے پاسے، جے $!A$، $\Th{\bar Q}$ وچ اے، تاں پہلے درجے دی
منطق وچ $!E \lif \lnot !A$ دا ثبوت اے۔ جے $\Th{T}$،
$!E \lif !A$ وی ثابت کرے، تاں $\Th{T}$، $\lnot !E$ ثابت
کردا اے۔ اوس صورت وچ $\Th{T}
\cup \Th{Q}$ متضاد ہووے گا۔ پر فرض مطابق
$\Th{T} \cup \Th{Q}$ بے تضاد اے۔ ایس لئی $\Th{T}$،
$!E \lif !A$ ثابت نہیں کردا تے $!A$، $C$ وچ نہیں۔

اسیں وکھایا اے کہ جے $!A$، $\Th{Q}$ وچ اے، تاں اوہ $C$
وچ اے۔ جے $!A$، $\Th{\bar Q}$ وچ اے، تاں اوہ
$\Complement{C}$ وچ اے۔ ایس طرح $C$ دونہاں سیٹاں نوں
قابلِ حساب طریقے نال وکھ کردا اے۔ ایہی اوہ تضاد اے جیہڑا
اسیں لبھ رہے ساں۔
\end{proof}

ایہ قضیہ بڑا مضبوط اے۔ مثال لئی، ایس توں ایہ نتیجہ نکلدا اے۔

\begin{cor}
حساب دی زبان دی پہلے درجے دی منطق (یعنی سیٹ
$\Setabs{!A}{\text{$!A$ پہلے درجے دی منطق وچ ثابت ہو سکدا اے}}$)
ناقابلِ فیصلہ اے۔
\end{cor}

\begin{proof}
پہلے درجے دی منطق $\emptyset$ دے منطقی نتیجیاں دا سیٹ اے۔
ایہ سیٹ~$\Th{Q}$ نال کٹھا بے تضاد اے۔
\end{proof}

\end{document}
